\section{Uniform cover counts for small second-norm residuals}
\label{qc-section}

We quantify the arithmetic covers of Section~\ref{pl-section} in a
fixed number field as the exponent varies. Write
\[
 F(a,b)=a^4+3a^3b+5a^2b^2+3ab^3+b^4=VQ^g,
 \qquad \gcd(a,b)=1,\quad a,b,V>0,\quad Q\ge7,
\]
with all displayed parameters integers. Put
$\lambda=\max\{1,\log V\}/g$. The results control the number of covers
and the heights of their coefficients. They supply no uniform bound
on the heights of points on those covers.

\subsection{The fixed field and the bad primes}

\begin{lemma}[A fixed splitting field]\label{qc-fixed-field}
Let $\zeta^2-\zeta+1=0$, $K_0=\mathbb Q(\zeta)$, and
$d_0=-1-3\zeta$. The splitting field of $F(x)=F(x,1)$ is
\[
 K=K_0(\sqrt{d_0},\sqrt{13}).
\]
It has degree eight, signature $(0,4)$, and unit rank three, and is
unramified outside $3$ and $13$.
\end{lemma}
\begin{proof}
With $t=x+x^{-1}$, one has $F(x)/x^2=t^2+3t+3$. The two values of
$t$ are $\zeta-2$ and $-1-\zeta$; the corresponding quadratic
discriminants are $d_0$ and $\bar d_0$, with product $13$. This gives
the stated splitting field. Neither $d_0$ nor $\bar d_0$ is a square
in $K_0$, since each has norm $13$. Also $\sqrt{13}\notin K_0$,
whose real elements are rational. For a nonsquare $d$ in a field of
characteristic zero, an element $t$ becomes a square in the quadratic
extension by $\sqrt d$ only if $t$ or $t/d$ was already a square:
expand $(u+v\sqrt d)^2=t$ and use $uv=0$. Applying this to
$t=13$, $d=d_0$, gives degree $2\cdot2\cdot2=8$.
The field contains the imaginary quadratic field $K_0$, so its
signature is $(0,4)$; the unit theorem gives rank three.

The discriminant of $F$ is $117$. At a prime outside $117$ its
residual polynomial is separable. Its roots lift in an unramified
extension after extending the residue field to split it, by simple-root
Hensel lifting. Hence its splitting field is unramified there.
The root differences are also units at all such primes.
\end{proof}

Fix roots $\alpha_1,\ldots,\alpha_4$ in $K$. Let $H$ be the class
number of $K$, and $w$ the number of its roots of unity. Their values
are not assumed. The standard ideal and unit inputs, including the full
logarithmic unit lattice, are supplied by
\cite[Chapters~3--5, Theorems~5.1 and~5.9]{MilneANT2020}.
There is no need to adjoin $g$-th roots of unity to the arithmetic
field: the geometric Kummer calculations take place over its algebraic
closure.

\begin{lemma}[Bad-prime depth]\label{qc-bad-depth}
For every primitive positive seed, $F(a,b)\equiv1\pmod3$, and
$13\mid F(a,b)$ implies $v_{13}(F(a,b))=1$.
\end{lemma}
\begin{proof}
Modulo three, $F(a,b)=(a^2+b^2)^2=1$. Modulo thirteen,
\[
 F(a,b)\equiv(a-b)^2(a^2+5ab+b^2).
\]
The quadratic factor has nonsquare discriminant $8$ modulo thirteen.
Thus $13\mid F$ forces $a\equiv b\not\equiv0\pmod{13}$. Writing
$a=b+13t$ and using $F(1,1)=13$, $F_x(1,1)=26$ gives
$F(b+13t,b)\equiv13b^4\pmod{169}$, proving exact depth one.
\end{proof}

\subsection{A count uniform in the exponent}

\begin{theorem}[Fixed-residual cover count]\label{qc-cover-count}
Let $g>8$ and let $V$ be a positive $g$-free integer. For fixed $g,V$,
all primitive positive seeds $F(a,b)=VQ^g$, with arbitrary $Q$, lift to
one of at most
\begin{equation}\label{qc-count}
 C_K4^{8\omega(V)}g^9,\qquad C_K=9^{32}H^4w^3,
\end{equation}
curves over $K$, of the form
\begin{equation}\label{qc-covers}
 y_i^g=\frac{x-\alpha_i}{\xi_i(x-\alpha_4)},\qquad i=1,2,3,
 \qquad \xi_i\in K^\times.
\end{equation}
Each smooth projective normalization is geometrically connected, with
degree $g^3$ over $\mathbb P^1$ and genus $1+g^2(g-2)$.
The constant is independent of $g,V,Q$; repetitions in this list are
allowed.
\end{theorem}
\begin{proof}
Set $L_i=a-\alpha_i b$ and factor ideals as
\[
 (L_i)=A_iB_i^g,
 \qquad 0\le v_{\mathfrak p}(A_i)<g.
\]
All these ideals are integral. Outside $3,13$, at most one $L_i$ is
divisible by a given prime $\mathfrak p$: otherwise the unit root
difference forces divisibility of both $a$ and $b$. Unramifiedness
and $\prod_iL_i=VQ^g$ show that the nonzero remainder valuation is
exactly $v_p(V)$. At each prime above $p\mid V$ it can be assigned
to one of four ideals $A_i$. There are at most eight primes above $p$,
giving at most $4^8$ choices per rational prime.

There is no contribution above three. At a prime above thirteen the
total valuation of $\prod_iL_i$ is at most its ramification index,
which is at most eight. As $g>8$, this entire contribution remains
in the $A_i$. Bound its four allocations by $9^4$ for each of at most
eight primes above thirteen. The possible ordered quadruples $(A_i)$
therefore number at most $9^{32}4^{8\omega(V)}$. These are deliberately
coarse bounds. The valuation argument also gives
\begin{equation}\label{qc-ideal-product}
 \prod_i A_i=(V),\qquad \operatorname N A_i\le V^8.
\end{equation}

For each $A_i$, the class $[B_i]$ solves
$[B_i]^g=[A_i]^{-1}$ and has at most $H$ possibilities. Choose an
integral ideal $C_i$ in each possible class and a generator $\delta_i$
of $A_iC_i^g$. Write $B_i=(t_i)C_i$; then
\[
 L_i=\delta_i\varepsilon_i t_i^g,
 \qquad \varepsilon_i\in\mathcal O_K^\times.
\]
The four ideal classes have at most $H^4$ choices. For the three
ratios $L_i/L_4$ only the three unit classes of
$\varepsilon_i/\varepsilon_4$ matter. The unit theorem gives
\[
 |\mathcal O_K^\times/(\mathcal O_K^\times)^g|
   =\gcd(g,w)g^3\le wg^3.
\]
Thus there are at most $w^3g^9$ ratio-class choices, proving
\eqref{qc-count}. These give actual $K$-points of
\eqref{qc-covers}, with $x=a/b$; positive $x$ avoids the roots.

Over the algebraic closure, valuation at each $\alpha_i$ for $i\le3$
proves the three Kummer classes independent modulo $g$, including
composite $g$. The degree is $g^3$. The inertia at each of the four
roots is $g$; at $\alpha_4$ it is the diagonal subgroup because the
three simple poles adjoin the same root of a uniformizer. Infinity
is unramified. Thus
\[
 2G_g-2=-2g^3+4g^2(g-1),\qquad G_g=1+g^2(g-2).
\]
\end{proof}

\begin{corollary}[Small residuals, fixed and varying]\label{qc-lambda-count}
If $g>8$ and $\lambda=\max\{1,\log V\}/g<\log2$, then $V$ is
$g$-free, and a fixed $(g,V)$ list has size $N_{g,V}$ satisfying
\begin{equation}\label{qc-fixed-logcount}
 \frac{\log\max\{1,N_{g,V}\}}g
 \le16\lambda+\frac{9\log g+\log C_K}g.
\end{equation}
For $0<L<\log2$, all seeds with exponent $g>8$ and
$1\le V\le\exp(Lg)$ are covered by a combined list of size at most
\begin{equation}\label{qc-union-count}
 C_Kg^9\exp(17Lg).
\end{equation}
In particular, when $L=L_g\to0$ and $g\to\infty$, the logarithm of
this combined list size divided by $g$ tends to zero.
\end{corollary}
\begin{proof}
For each $p$, $v_p(V)\le\log V/\log2<g$. Also
$\omega(V)\le\log V/\log2$, so
$8\log4\,\omega(V)\le16\log V$. This proves
\eqref{qc-fixed-logcount}. For the union, there are at most
$\exp(Lg)$ integer residuals, each with at most
$C_Kg^9\exp(16Lg)$ covers. Sum these bounds.
\end{proof}

The factor for the choice of $V$ is retained in
\eqref{qc-union-count}; a bound for one fixed residual is not silently
used as a bound for all moving residuals. The construction applies
to prime exponents, twice a prime, and other $g>8$ alike.

\subsection{Coefficient heights and a strict unit-class boundary}

Write $h_K$ for the absolute logarithmic Weil height. For a reduced
positive rational $a/b$, it is $\log\max(a,b)$.

\begin{theorem}[Coefficient normalization]\label{qc-coefficient-height}
The lists above may be chosen so that
\begin{equation}\label{qc-xi-height}
 h_K(\xi_i)\le2\log V+C'_Kg.
\end{equation}
Consequently the coefficient heights of the three equations are
$O_K(\log V+g)$, uniformly in $g,V,Q$.
\end{theorem}
\begin{proof}
Choose a fixed integral representative of every ideal class, of norm
at most $M_K$. By \eqref{qc-ideal-product},
\[
 \log\operatorname N(A_iC_i^g)\le8\log V+g\log M_K.
\]
A principal integral ideal $J$ has a generator $\delta$ satisfying
\begin{equation}\label{qc-balanced-generator}
 h_K(\delta)\le\tfrac18\log\operatorname N J+B_K.
\end{equation}
Indeed, subtract the average $\log\operatorname N J/8$ from the
four complex logarithmic absolute values of an integral generator.
Reduce this sum-zero vector modulo the full logarithmic unit lattice
into a fixed bounded fundamental domain. Multiplication by the
corresponding unit keeps the generator integral and bounds every
logarithmic absolute value by that average plus a fixed constant.
The average is nonnegative; summing positive parts proves
\eqref{qc-balanced-generator}.

Choose the $\delta_i$ in this way. Then
\[
 h_K(\delta_i)\le\log V+\tfrac g8\log M_K+B_K.
\]
For fixed fundamental units $\epsilon_1,\epsilon_2,\epsilon_3$,
every unit class modulo $g$-th powers has a representative
$\rho=\omega\prod_j\epsilon_j^{e_j}$, with $0\le e_j<g$ and
$\omega$ a root of unity. Its height is at most
$g\sum_jh_K(\epsilon_j)$. Taking
$\xi_i=(\delta_i/\delta_4)\rho$ for each required ratio class and
using the multiplication and inversion height inequalities gives
\eqref{qc-xi-height}. These choices preserve every actual lift.
\end{proof}

\begin{lemma}[Linear height is necessary for some unit classes]
\label{qc-unit-lower-bound}
For a fixed fundamental unit $\epsilon_1$ there is $c_K>0$ such that,
for every $g\ge3$, every representative in $K^\times$ of the class
$[\epsilon_1^{\lfloor g/2\rfloor}]$ modulo $K^{\times g}$ has
height at least $c_Kg$.
\end{lemma}
\begin{proof}
Write $k=\lfloor g/2\rfloor$ and $\xi=\epsilon_1^kt^g$. If $t$
is not a unit, some finite valuation of $\xi$ has absolute value at
least $g$. Applying the denominator contribution to $\xi$ or to
$\xi^{-1}$, whose heights agree, gives
$h_K(\xi)\ge g\log2/8$.

If $t$ is a unit, its coordinates in the fundamental-unit basis give
the exponent vector $(k+gn_1,gn_2,gn_3)$ for $\xi$, up to torsion.
Its first coordinate has absolute value at least $\lfloor g/2\rfloor
\ge g/3$. Weil height on the real logarithmic unit space is a norm:
the weighted sum is zero, so the sum of positive parts is half the
weighted absolute sum. The fundamental-unit logarithms are linearly
independent. Equivalence with the maximum norm in these fixed coordinates
therefore gives $h_K(\xi)\ge c'_Kg/3$.
\end{proof}

Thus arbitrary unit classes cannot all be normalized to height $o(g)$.
They are already present in the unrestricted class list with trivial
residual ideals. This does not show that these particular classes have
actual positive seeds. To remove their possible contribution one must
use the compatibility of actual points, rather than merely choosing
different class representatives. The order-$g$ term in
\eqref{qc-xi-height} is not absorbed into $O(\log V)$.

\subsection{The remaining point-height problem}

Put $u=ab/(a+b)^2\in(0,1/4]$. Then
\[
 \frac{F(a,b)}{(a+b)^4}=1-u+u^2,\qquad
 \frac{13}{16}(a+b)^4\le F(a,b)<(a+b)^4.
\]
Every actual seed therefore satisfies
\begin{equation}\label{qc-point-height-lower}
 h_K(a/b)\ge\frac{g\log Q+\log V}{4}-\log2
          \ge\frac{g\log7}{4}-\log2.
\end{equation}
A family with $\lambda\to0$ would require points of height at least
linear in $g$, on a list whose logarithmic size is $o(g)$ and whose
coefficient heights are $O(g)$. These facts are compatible: the
number of possible curves gives no height bound for a point on one
curve. No estimate uniform in $g$ for these point heights has been
proved here. Moving residual support, large-prime exponents, and the
two-step small-residual threshold remain unresolved.

This section proves ordinary arithmetic-geometric reductions. It does
not assert a uniform Faltings theorem or Lean verification of the
number-field, unit-lattice, or geometric arguments.
